\documentclass[12pt]{article}

% Packages
\usepackage[margin=1in]{geometry}
\usepackage{titlesec}
\usepackage{titletoc}
\usepackage{hyperref}
\usepackage{longtable}
\usepackage{booktabs}
\usepackage{array}
\usepackage{xcolor}
\usepackage{fancyhdr}
\usepackage{graphicx}
\usepackage{parskip}
\usepackage{enumitem}
\usepackage{tabularx}

% Hyperlink setup
\hypersetup{
    colorlinks=true,
    linkcolor=black,
    urlcolor=blue,
    citecolor=black
}

% Header / Footer
\pagestyle{fancy}
\fancyhf{}
\rhead{CS3338 -- Group 07}
\lhead{Software Design Document}
\rfoot{Page \thepage}
\lfoot{Automated Vertical Basketball Highlight System}

% Section formatting
\titleformat{\section}{\large\bfseries}{}{0em}{}[\titlerule]
\titleformat{\subsection}{\normalsize\bfseries}{}{0em}{}

% ─────────────────────────────────────────────
\begin{document}

% ══════════════════════════════════════════════
%  COVER PAGE
% ══════════════════════════════════════════════
\begin{titlepage}
    \centering
    \vspace*{2cm}

    {\LARGE\bfseries Software Design Document\\[0.5em]}
    {\Large Automated Vertical Basketball Highlight System\\[2em]}

    \rule{\linewidth}{0.5pt}\\[1em]

    {\large CS3338 -- Software Engineering\\
    California State University, Los Angeles\\[1em]}

    {\large Group 07\\[2em]}
    {\normalsize Christopher Ayala, Edwin Hui, Kazuho Igarashi, Nestor Adame Delgado, Ovi Ahmed\\[2em] }

    \begin{tabular}{ll}
        \textbf{Version:}  & 1.0                   \\
        \textbf{Date:}     & \today                \\
        \textbf{Status:}   & Snapshot 1 -- Initial Draft \\
    \end{tabular}

    \vfill
    {\small Jira Project: \href{https://cs3338-group-07.atlassian.net}{https://cs3338-group-07.atlassian.net}}
\end{titlepage}

% ══════════════════════════════════════════════
%  TABLE OF CONTENTS
% ══════════════════════════════════════════════
\tableofcontents
\newpage

% ══════════════════════════════════════════════
%  VERSIONS TABLE
% ══════════════════════════════════════════════
\section*{Version History}
\addcontentsline{toc}{section}{Version History}

\begin{tabularx}{\textwidth}{|c|c|X|c|}
    \hline
    \textbf{Version} & \textbf{Date} & \textbf{Description} & \textbf{Snapshot} \\
    \hline
    1.0 & \today & Initial draft. System architecture, UI design, and Snapshot 1 objectives defined. & Snapshot 1 \\
    \hline
    1.1 & TBD & [Snapshot 2 -- To be updated] & Snapshot 2 \\
    \hline
    1.2 & TBD & [Snapshot 3 -- To be updated] & Snapshot 3 \\
    \hline
    1.3 & TBD & [Snapshot 4 -- To be updated] & Snapshot 4 \\
    \hline
\end{tabularx}

\newpage

% ══════════════════════════════════════════════
%  SECTION 1 -- INTRODUCTION
% ══════════════════════════════════════════════
\section{Introduction}

\subsection{Purpose}
This Software Design Document (SDD) describes the architectural and detailed design
of the Automated Vertical Basketball Highlight System. It serves as a technical
reference for the development team throughout all project snapshots, providing a
clear blueprint of system components, data flow, and user interface decisions.

\subsection{Intended Audience}
\begin{itemize}[noitemsep]
    \item Development team members (Group 07)
    \item Course instructor and teaching assistants
    \item Future maintainers or contributors to the project
\end{itemize}

\subsection{Project Overview}
The Automated Vertical Basketball Highlight System is a Python-based application
that processes standard 16:9 basketball game footage and automatically generates
9:16 vertical highlight clips optimized for mobile platforms such as TikTok,
Instagram Reels, and YouTube Shorts.

The system uses YOLOv8 for real-time object detection (ball, players, hoop) and a
Kalman filter for smooth subject tracking. A Flask web application provides the
user interface, and the entire stack is containerized with Docker for portability
and ease of deployment.

\subsection{Scope}
\begin{itemize}[noitemsep]
    \item Accepts uploaded or locally stored basketball game video files.
    \item Detects and tracks the ball, players, and hoop using computer vision.
    \item Automatically crops and reframes footage from 16:9 to 9:16 format.
    \item Identifies highlight-worthy moments (e.g., shots, dunks, fast breaks).
    \item Exports final highlight clips via the Flask UI.
\end{itemize}

\subsection{Document Overview}
The remainder of this document is organized as follows:
\begin{itemize}[noitemsep]
    \item \textbf{Section 2} -- System Architecture
    \item \textbf{Section 3} -- User Interface Design
    \item \textbf{Section 4} -- Snapshot Revision Notes
    \item \textbf{Section 5} -- Glossary
    \item \textbf{Section 6} -- References
\end{itemize}

\newpage

% ══════════════════════════════════════════════
%  SECTION 2 -- SYSTEM ARCHITECTURE
% ══════════════════════════════════════════════
\section{System Architecture}

\subsection{High-Level Architecture}
The system is divided into four primary layers:

\begin{enumerate}[noitemsep]
    \item \textbf{Presentation Layer} -- Flask web UI served to the user's browser.
    \item \textbf{Application Layer} -- Python business logic coordinating detection,
          tracking, and clip extraction.
    \item \textbf{Computer Vision Layer} -- YOLOv8 model for object detection;
          Kalman filter for trajectory smoothing.
    \item \textbf{Storage Layer} -- Local filesystem for input videos and output clips;
          optionally extensible to cloud storage.
\end{enumerate}

\subsection{Component Breakdown}

\begin{tabularx}{\textwidth}{|l|X|}
    \hline
    \textbf{Component} & \textbf{Description} \\
    \hline
    \texttt{Flask UI} & Web interface for video upload, processing status, and
                        clip download. Runs on port 5001 to avoid macOS AirPlay conflict. \\
    \hline
    \texttt{YOLO Detector} & YOLOv8 model (Ultralytics) fine-tuned or applied
                              zero-shot to detect basketball, players, and hoop. \\
    \hline
    \texttt{Kalman Tracker} & Kalman filter applied per-frame to smooth
                               bounding-box trajectories and reduce jitter. \\
    \hline
    \texttt{Crop Engine} & Computes the optimal 9:16 crop window per frame
                           based on tracked objects and reframes the video. \\
    \hline
    \texttt{Highlight Scorer} & Heuristic or ML-based module that scores segments
                                 and selects highlight clips. \\
    \hline
    \texttt{Export Module} & Uses FFmpeg to encode and export the final 9:16 MP4. \\
    \hline
\end{tabularx}

\subsection{Data Flow}
\begin{enumerate}[noitemsep]
    \item User uploads a video via the Flask UI.
    \item The application layer invokes the YOLO Detector frame-by-frame.
    \item Detection results are passed to the Kalman Tracker to produce smooth trajectories.
    \item The Crop Engine uses trajectories to calculate the 9:16 crop window per frame.
    \item The Highlight Scorer evaluates segments and selects clip boundaries.
    \item The Export Module encodes the selected segments and saves the output clip.
    \item The Flask UI provides a download link to the user.
\end{enumerate}

\subsection{Technology Stack}

\begin{tabularx}{\textwidth}{|l|l|X|}
    \hline
    \textbf{Layer} & \textbf{Technology} & \textbf{Purpose} \\
    \hline
    Web Framework  & Flask (Python)      & UI and API routing \\
    \hline
    Object Detection & YOLOv8 (Ultralytics) & Detect ball, players, hoop \\
    \hline
    Tracking       & Kalman Filter (NumPy/SciPy) & Smooth object trajectories \\
    \hline
    Video Processing & OpenCV, FFmpeg    & Frame extraction and encoding \\
    \hline
    Containerization & Docker, docker-compose & Deployment and environment isolation \\
    \hline
    Project Management & Jira (Atlassian) & Sprint planning and task tracking \\
    \hline
    Version Control & GitHub            & Source code and document storage \\
    \hline
    Testing        & TestRail           & Test case management \\
    \hline
\end{tabularx}

\subsection{Docker Services}
The application is containerized using \texttt{docker-compose}. The following
services are defined:

\begin{itemize}[noitemsep]
    \item \textbf{web} -- Flask application container (Python 3.11, port 5001).
    \item \textbf{worker} -- Background processing container for CV pipeline.
    \item \textbf{ffmpeg} -- Sidecar container or shared binary for video encoding.
\end{itemize}

\newpage

% ══════════════════════════════════════════════
%  SECTION 3 -- USER INTERFACE
% ══════════════════════════════════════════════
\section{User Interface}

\subsection{Overview}
The user interface is a browser-based web application built with Flask and
rendered using Jinja2 templates with Bootstrap for responsive styling.
No frontend JavaScript framework is used, keeping the stack lightweight.

\subsection{Pages and Views}

\begin{tabularx}{\textwidth}{|l|X|}
    \hline
    \textbf{Page / Route} & \textbf{Description} \\
    \hline
    \texttt{/} (Home)     & Landing page with project description and upload button. \\
    \hline
    \texttt{/upload}      & File upload form; accepts MP4, MOV, AVI video files. \\
    \hline
    \texttt{/process}     & Processing status page showing real-time progress. \\
    \hline
    \texttt{/result}      & Download page displaying the generated highlight clip. \\
    \hline
    \texttt{/about}       & Project information and team credits. \\
    \hline
\end{tabularx}

\subsection{User Interaction Flow}
\begin{enumerate}[noitemsep]
    \item User navigates to the home page.
    \item User uploads a basketball game video via the upload form.
    \item The server begins processing; the user is redirected to the status page.
    \item Once complete, the user is redirected to the result page.
    \item The user downloads the 9:16 highlight clip.
\end{enumerate}

\subsection{Database / Storage}
In Snapshot 1, no persistent database is used. Video files and output clips are
stored on the local filesystem within Docker volumes. Future snapshots may
introduce a lightweight database (e.g., SQLite or PostgreSQL) to store processing
history and user sessions.

\newpage

% ══════════════════════════════════════════════
%  SECTION 4 -- SNAPSHOT REVISION NOTES
% ══════════════════════════════════════════════
\section{Snapshot Revision Notes}

\subsection{Snapshot 1 -- Initial Design (Current)}
\begin{itemize}[noitemsep]
    \item Defined system architecture and component responsibilities.
    \item Established technology stack and Docker service structure.
    \item Designed Flask UI page layout and user interaction flow.
    \item Set up Jira project and initial sprint backlog.
    \item Configured GitHub repository structure.
\end{itemize}

\subsection{Snapshot 2 -- Checkpoint 1}
At this checkpoint, the core computer vision pipeline was integrated and validated
end-to-end. The following design updates were made:

\begin{itemize}[noitemsep]
    \item \textbf{Highlight Scorer module introduced.} A heuristic scoring module
          was added to the application layer. It analyzes per-frame detection
          confidence and ball proximity to the hoop to identify candidate highlight
          segments. Clips scoring above a configurable threshold are flagged for export.
    \item \textbf{Processing status via polling.} The \texttt{/process} route was
          updated to support lightweight AJAX polling so the browser refreshes
          progress without a full page reload.
    \item \textbf{Volume mount added for output clips.} The \texttt{docker-compose}
          configuration was updated to mount a named volume (\texttt{clips\_output})
          shared between the \texttt{web} and \texttt{worker} containers, ensuring
          processed files persist across container restarts.
    \item \textbf{New dependency: \texttt{filterpy}.} The Kalman filter implementation
          was migrated from a custom NumPy implementation to the \texttt{filterpy}
          library for improved maintainability and accuracy.
    \item \textbf{TestRail test cases added.} Test cases covering video upload,
          detection accuracy, and clip export were defined and executed in TestRail.
\end{itemize}

\subsection{Snapshot 3 -- Checkpoint 2}
This checkpoint focused on improving output quality and expanding the user interface.
The following design updates were made:

\begin{itemize}[noitemsep]
    \item \textbf{Multi-clip export.} The Export Module was redesigned to support
          exporting multiple highlight clips from a single input video. Users can
          now preview a list of detected highlights and selectively download clips
          via the \texttt{/result} page.
    \item \textbf{Clip preview thumbnails.} The Flask UI was updated to generate
          a thumbnail image for each highlight clip using FFmpeg's
          \texttt{-vframes 1} flag. Thumbnails are displayed on the result page
          before download.
    \item \textbf{Configurable crop padding.} A padding parameter was introduced
          in the Crop Engine to add a configurable margin around the tracked
          subject, reducing aggressive cropping on fast-moving plays.
    \item \textbf{New route: \texttt{/history}.} A session-based history page was
          added, listing previously processed videos and their output clips within
          the current browser session.
    \item \textbf{New dependencies: \texttt{Pillow}, \texttt{flask-session}.}
          Pillow is used for thumbnail generation; flask-session provides
          server-side session storage for the history feature.
    \item \textbf{TestRail test cases updated.} Additional test cases covering
          multi-clip selection, thumbnail rendering, and session history were
          added and executed.
\end{itemize}

\subsection{Snapshot 4 -- Final}
The final snapshot addressed polish, stability, and documentation. The following
changes were made:

\begin{itemize}[noitemsep]
    \item \textbf{Error handling improvements.} Graceful error pages were added
          for unsupported file formats, processing failures, and timeout events.
          All exceptions in the worker container are logged to a file mounted
          via Docker volume.
    \item \textbf{Processing time optimization.} Frame sampling was introduced
          in the YOLO Detector: only every other frame is passed through the
          model during the initial highlight-scoring pass, reducing total
          processing time by approximately 40\% with minimal accuracy loss.
    \item \textbf{UI refinement.} Bootstrap layout was updated for improved
          mobile responsiveness. The upload form now displays file size
          validation feedback before submission.
    \item \textbf{Docker Compose finalized.} The \texttt{docker-compose.yml}
          was updated with health checks for the \texttt{web} and \texttt{worker}
          services and a restart policy of \texttt{unless-stopped}.
    \item \textbf{Future Work.} The following features were identified but not
          implemented within the project timeline:
          \begin{itemize}[noitemsep]
              \item Cloud storage integration (AWS S3) for persistent clip hosting.
              \item Fine-tuned YOLOv8 model trained on a labeled basketball dataset
                    for improved detection accuracy in low-light and crowded scenes.
              \item User authentication and per-user clip libraries.
              \item Real-time streaming input via RTMP for live game processing.
          \end{itemize}
\end{itemize}

\newpage

% ══════════════════════════════════════════════
%  SECTION 5 -- GLOSSARY
% ══════════════════════════════════════════════
\section{Glossary}

\begin{tabularx}{\textwidth}{|l|X|}
    \hline
    \textbf{Term / Acronym} & \textbf{Definition} \\
    \hline
    API   & Application Programming Interface \\
    \hline
    CV    & Computer Vision \\
    \hline
    Docker & Containerization platform for packaging and deploying applications \\
    \hline
    FFmpeg & Open-source multimedia framework for video encoding and processing \\
    \hline
    Flask & Lightweight Python web framework \\
    \hline
    FPS   & Frames Per Second \\
    \hline
    Jira  & Project management and issue tracking tool by Atlassian \\
    \hline
    Kalman Filter & Algorithm for estimating the state of a linear dynamic system from noisy measurements \\
    \hline
    MP4   & MPEG-4 video container format \\
    \hline
    SDD   & Software Design Document \\
    \hline
    SRS   & Software Requirements Specification \\
    \hline
    UI    & User Interface \\
    \hline
    UX    & User Experience \\
    \hline
    YOLO  & You Only Look Once -- real-time object detection algorithm \\
    \hline
    YOLOv8 & Version 8 of the YOLO object detection model by Ultralytics \\
    \hline
    9:16  & Vertical aspect ratio used for mobile-first video platforms \\
    \hline
    16:9  & Standard widescreen aspect ratio for video footage \\
    \hline
\end{tabularx}

\newpage

% ══════════════════════════════════════════════
%  SECTION 6 -- REFERENCES
% ══════════════════════════════════════════════
\section{References}

\begin{enumerate}[noitemsep]
    \item Ultralytics YOLOv8 Documentation. \url{https://docs.ultralytics.com}
    \item Flask Documentation. \url{https://flask.palletsprojects.com}
    \item OpenCV Documentation. \url{https://docs.opencv.org}
    \item FFmpeg Documentation. \url{https://ffmpeg.org/documentation.html}
    \item Docker Documentation. \url{https://docs.docker.com}
    \item Atlassian Jira -- Group 07 Project. \url{https://cs3338-group-07.atlassian.net}
    \item Welch, G. \& Bishop, G. (1995). \textit{An Introduction to the Kalman Filter}.
          University of North Carolina at Chapel Hill.
    \item Redmon, J. et al. (2016). \textit{You Only Look Once: Unified, Real-Time
          Object Detection}. CVPR 2016.
\end{enumerate}

\end{document}